

\section{Lower bound on PC degree for Tensor Isomorphism from Graph Isomorphism} \label{sec:lowerboundGI}
\begin{definition} \label{def:GI}
Given two graphs $G,H$ with adjacency matrices $A,B$ (resp.), the equations for \textsc{Graph Isomorphism} (the same as those used by Berkholz \& Grohe \cite{BerkholzG15,BerkholzG16}) are as follows. Let $Z$ be an $n \times n$ matrix of variables $z_{ij}$ (where the intended interpretation is that $z_{ij}=1$ iff an isomorphism maps vertex $i \in V(G)$ to vertex $j \in V(H)$). We say that a partial map, which sends $(i,i') \mapsto (j,j')$ is a local isomorphism if (1) $i=i'$ iff $j=j'$ (it's a well-defined map) and (2) $(i,i') \in E(G) \Leftrightarrow (j,j') \in E(H)$. (One may also do \textsc{Colored Graph Isomorphism} and require that the colors match, $c(i) = c(j), c(i')=c(j')$.) Then the equations are:
\[
\begin{array}{rll}
z_{ij}^2 - z_{ij} & \forall i,j & \text{All variables $\{0,1\}$-valued} \\
1 - \sum_i z_{ij} & \forall j & \text{each $j \in V(H)$ is mapped to from exactly one vertex} \\
1 - \sum_j z_{ij} & \forall i & \text{each $i \in V(G)$ maps to exactly one vertex} \\
z_{ij} z_{i'j'} & & \text{Whenever $(i,i') \mapsto (j,j')$ is not a local isomorphism}.
\end{array}
\]
\end{definition}

In this section, we prove a lower bound on PC (and SoS) for \textsc{TI}, by reducing from \textsc{GI} and using the known lower bounds on \textsc{GI} \cite{BerkholzG15,BerkholzG16}. Specifically, we show

\begin{theorem}
Over any field, there are instances of \textsc{Tensor Isomorphism} of size $O(n) \times O(n) \times O(n)$ that require PC degree $\Omega(n)$ to refute. The same holds over the reals for SoS degree.
\end{theorem}

\begin{proof}
Berkholz and Grohe \cite{BerkholzG15,BerkholzG16} show the same statement for $n$-vertex graphs of bounded vertex degrees, with the same PC/SoS degree bound. In Proposition~\ref{prop:GIMon} we show that \textsc{GI} reduces to \textsc{Monomial Code Equivalence} by a (2,4)-many-one reduction that turns $n$-vertex, $m$-edge graphs into $m \times (3m+n)$ matrices. in Proposition~\ref{prop:MonTI} we show that \textsc{Monomial Code Equivalence} reduces to \textsc{TI} by a (2,4)-many-one reduction that turns $k \times N$ matrices into $(k + 2N) \times N \times (1 + 2N)$ tensors. By Lemma~\ref{lem:PCred}, this completes the proof.
\end{proof}



To reduce from \textsc{GI} to \textsc{TI} we use the following intermediate problem. A matrix is \emph{monomial} if it has exactly one nonzero entry in each row and column; equivalently, a monomial matrix is the product of a permutation matrix and an invertible diagonal matrix. 

\begin{definition}
\textsc{Monomial Code Equivalence} is the problem: given two $k \times n$ matrices $C,C'$, do there exist matrices $X,Y$ such that $XCY^T = C'$ where $X$ is invertible and $Y$ is invertible and monomial? Given two such matrices $C,C'$, the equations for \textsc{Monomial Code Equivalence} are as follows. There are $2(k^2 + n^2)$ variables arranged into matrices $X,X'$ (of size $k \times k$) and $Y,Y'$ (of size $n \times n$). The equations are
\[
XCY^T = C' \qquad XX' = X'X =\Id \qquad YY' = Y'Y = \Id
\]
and
\begin{align*}
& y_{ij}y_{ij'} (\forall i \forall j \neq j') \qquad y_{ij} y_{i'j} (\forall i \neq i', \forall j) \\
& y'_{ij}y'_{ij'} (\forall i \forall j \neq j') \qquad y'_{ij} y'_{i'j} (\forall i \neq i', \forall j) 
\end{align*}
(Note: there are no equations forcing the variables to take on values in $\{0,1\}$.)
\end{definition}


\begin{proposition} \label{prop:GIMon}
The reduction of Petrank \& Roth \cite{PR} from \textsc{Graph Isomorphism} to \textsc{Linear Code Equivalence} over $\F_2$ in fact gives a (2,4)-many-one reduction from \textsc{Graph Isomorphism} to \textsc{Monomial Code Equivalence} (sic!) over any field.
\end{proposition}

\begin{proof}
The reduction of Petrank \& Roth is as follows: given a simple undirected graph $G$ with $n$ vertices and $m$ edges, let $D(G)$ be its $m \times n$ incidence matrix: $D_{e,v} = 1$ iff $v \in e$ and is 0 otherwise, and let $M(G)$ be the $m \times (3m + n)$ matrix 
\[
M(G) = \left[
\begin{array}{c|c|c|c}
I_m & I_m & I_m & D(G)
\end{array}
\right].
\]

\textbf{Many-one reduction.} It was previously shown (over $\F_2$ in \cite{PR} and over arbitrary fields in \cite[Lem.~II.4]{GrochowLie}) that this gives a many-one reduction to \textsc{Permutational Code Equivalence}. Here we observe that the same reduction also gives a reduction to \textsc{Monomial Code Equivalence}. Thus, all that remains to show is that if $M(G)$ and $M(H)$ are monomially equivalent, then $G$ must be isomorphic to $H$. 

In fact, what was shown in \cite{PR} (over arbitrary fields in \cite{GrochowLie}) is that, up to permutation and scaling of the rows, $M(G)$ is the unique generator matrix of its code satisfying the following properties: (1) $M(G)$ is $m \times (3m+n)$, (2) each row has Hamming weight $\leq 5$, (3) any linear combination that includes two or more rows with nonzero coefficients has Hamming weight $\geq 6$. 

Now, suppose $(X,Y)$ is a monomial equivalence of the codes $M(G), M(H)$. Then the rowspans of $M(G) Y^T$ and $M(H)$ are the same. Since $Y$ is monomial, if we consider just the supports of the rows of $M(G) Y^T$, up to re-ordering the rows, by the preceding paragraph, those supports must be the same as the supports of the rows of $M(H)$. Thus $X$ must also be monomial. Say $X = DP$ and $Y=EQ$ where $D,E$ are diagonal and $P,Q$ are permutation matrices. Then $P M(G) Q^T$ has the same support as $X M(G) Y^T = M(H)$, and since $P$ and $Q$ are permutation matrices and $M(G)$ and $M(H)$ have all entries in $\{0,1\}$, we must have $P M(G) Q^T = M(H)$. Thus $M(G)$ and $M(H)$ are in fact equivalent by a \emph{permutation} matrix (in place of the monomial matrix $Y$). Thus, by the fact that $(G,H) \mapsto (M(G), M(H))$ was a reduction to \textsc{Permutational Code Equivalence}, we conclude that $G \cong H$.


\textbf{Low-degree PC reduction.} Let $X,X',Y,Y'$ be the variable matrices in the equations for \textsc{Monomial Code Equivalence} of $M(G), M(H)$, and let $Z$ be the variable matrix in the equations for \textsc{Graph Isomorphism} of $G,H$. Let $n = |V(G)|, m = |E(G)|$; so, $X,X'$ are of size $m$, $Y,Y'$ are of size $3m+n$, and $Z$ is of size $n$. 

Let $Z^{(2)}$ denote the $\binom{n}{2} \times \binom{n}{2}$ matrix whose $(\{i,i'\}, \{j,j'\})$ entry is $z_{ij} z_{i'j'} + z_{ij'} z_{i'j}$. The idea is that if $Z$ is a map on the vertices, then $Z^{(2)}$ is the corresponding map on the edges; the two terms come from the fact that the edge $\{i,i'\}$ can be mapped to the edge $\{j,j'\}$ either by $(i,i') \mapsto (j,j')$ or by $(i,i') \mapsto (j',j)$. Note that, since $Z$ is a permutation matrix, at most one of these terms is nonzero, and thus $Z^{(2)}$ is also a $\{0,1\}$-matrix (in fact, a permutation matrix). Let $Z^{(2)}_E$ denote the $|E| \times |E|$ submatrix of $Z^{(2)}$ all of whose row indices are $\{i,i'\} \in E(G)$ and all of whose column indices are $\{j,j'\} \in E(H)$. Note also that $(Z^{(2)}_E)^T = (Z^T)^{(2)}_E$, so we use these notations interchangeably for convenience.

Now consider the following substitution:
\[
\begin{array}{ll}
X \mapsto (Z^{(2)}_E)^T & Y \mapsto (Z^T)^{(2)}_E \oplus (Z^T)^{(2)}_E \oplus (Z^T)^{(2)}_E \oplus (Z^T)  \\
X' \mapsto Z^{(2)}_E & Y' \mapsto Z^{(2)}_E \oplus Z^{(2)}_E \oplus Z^{(2)}_E \oplus Z \\
\end{array}
\]
After making these substitutions in the equations for \textsc{Monomial Code Equivalence} of $M(G), M(H)$, we get the equations
\begin{equation} \label{eq:monGI}
(Z^{(2)}_E)^T Z^{(2)}_E = Z^{(2)}_E (Z^{(2)}_E)^T = \Id_m \qquad (Z^{(2)}_E)^T D(G) Z = D(H) \qquad ZZ^T = Z^T Z = \Id_n
\end{equation}
along with equations saying that $Z$ and $Z^{(2)}_E$ are monomial.

We now show how to derive these equations in low-degree PC from the \textsc{GI} equations. 

The monomial equations for $Z$ are part of the \textsc{GI} equations, so there is nothing to do for those.

The monomial equations for $Z^{(2)}_E$ are of the form $(z_{ij} z_{i'j'} + z_{ij'} z_{i'j})(z_{k\ell} z_{k'\ell'} + z_{k\ell'} z_{k'\ell})$ where either (1) $\{i,i'\} = \{k,k'\}$ and $\{j,j'\} \neq \{\ell,\ell'\}$ or (2) vice versa. We expand out to get 
\[
z_{ij} z_{i'j'} z_{k\ell} z_{k'\ell'} + z_{ij} z_{i'j'} z_{k\ell'} z_{k'\ell} + z_{ij'} z_{i'j}z_{k\ell} z_{k'\ell'} + z_{ij'} z_{i'j} z_{k\ell'} z_{k'\ell}
\]
We show how to get this equation in case (1); case (2) follows similarly, \emph{mutatis mutandis}. In case (1), without loss of generality suppose that $i=k$, $i'=k'$, and $j \notin \{\ell,\ell'\}$. The first two terms are divisible by the \textsc{GI} equations $z_{ij} z_{i \ell}$ (since $i=k$ and $j \neq \ell$), the third term is divisible by $z_{i'j} z_{i' \ell'}$ (since $i'=k'$ and $j \neq \ell'$), and the last term is divisible by $z_{i'j} z_{i' \ell}$ similarly.

Next, the equations $ZZ^T=\Id_n$ are, expanded out, 
\[
\sum_j z_{ij} z_{ij} -1 (\forall i) \qquad \sum_j z_{ij} z_{kj} (\forall i \neq k).
\]
The first is gotten by linear combination from $1-\sum_j z_{ij}$ and the Boolean axioms $z_{ij}^2 - z_{ij}$. The second is a linear combination of the monomial axioms $z_{ij} z_{kj}$ (part of the local non-isomorphism axioms). Similarly for $Z^T Z = \Id$, using $1 - \sum_i z_{ij}$ instead.

Next, we expand out the equations $Z^{(2)}_E (Z^T)^{(2)}_E = \Id_m$, to get\footnote{We use the notation $\sum_{\{j,j'\} \in E(H)}$ to denote a sum in the index of summation takes on the value $e \in E(H)$ for each edge of $H$ exactly once. Because our edges are undirected, we only use such sums when the summand expression is itself invariant under swapping the roles of $j,j'$. If so desired, one could equivalently say $\sum_{j < j', \{j,j'\} \in E(H)}$.}
\[
\sum_{\{j,j'\} \in E(H)} (z_{ij} z_{i'j'} + z_{ij'} z_{i'j})(z_{kj} z_{k'j'} + z_{k'j} z_{kj'}) - \delta_{\{i,i'\}, \{k,k'\}} (\forall \{i,i'\}, \{k,k'\} \in E(G))
\]
Thus, for $\{i,i'\} \neq \{k,k'\}$, we need to derive
\[
\sum_{\{j,j'\} \in E(H)} \left( z_{ij} z_{i'j'} z_{kj} z_{k'j'}+ z_{ij'} z_{i'j} z_{kj} z_{k'j'} + z_{ij} z_{i'j'} z_{k'j} z_{kj'} + z_{ij'} z_{i'j} z_{k'j} z_{kj'} \right).
\]
Without loss of generality, suppose that $i \notin \{k,k'\}$. Then the first two terms of each summand are divisible by the \textsc{GI} equation $z_{ij} z_{kj}$, the third term is divisible by $z_{ij} z_{k'j}$, and the last term is divisible by $z_{ij'} z_{kj'}$. 
On the other hand, when $\{i,i'\} = \{k,k'\}$, we need to derive
\[
-1 + \sum_{\{j,j'\} \in E(H)} \left( z_{ij}^2 z_{i'j'}^2 + 2 z_{ij'} z_{i'j} z_{ij} z_{i'j'} + z_{ij'}^2 z_{i'j}^2 \right).
\]
The middle terms of each summand are divisible by the \textsc{GI} equations $z_{ij'} z_{ij}$. For the first and third terms, we can use the Boolean axioms to remove the squares, and thus we are left to derive
\begin{equation} \label{eq:Z2}
-1 + \sum_{\{j,j'\} \in E(H)} \left(z_{ij} z_{i'j'} + z_{ij'} z_{i'j} \right)
\end{equation}
We derive this from the \textsc{GI} equations as follows. Consider $(\sum_{j} z_{ij} -1)(\sum_{j'} z_{i'j'} -1) + (\sum_{j} z_{ij} -1) + (\sum_{j'} z_{i'j'} -1)$ and break up the resulting sum according to whether $j = j'$, $\{j,j'\} \in E(H)$ or $\{j,j'\} \notin E(H)$. Then we get
\[
\sum_j z_{ij} z_{i'j}  + \sum_{j,j' : \{j,j'\} \in E(H)} z_{ij} z_{i'j'} + \sum_{j \neq j' \{j,j'\} \notin E(H)} z_{ij} z_{i'j'} - 1
\]
Every summand in the first sum is a monomial axiom since $i \neq i'$. Every summand in the third sum is a local non-isomorphism axiom, since $\{i,i'\} \in E(G)$ but $\{j,j'\} \notin E(H)$. Note that every edge $\{j,j'\}$ of $E(H)$ is represented twice in the middle sum: once as $(j,j')$ and once as $(j',j)$. Thus, the above simplifies to 
\[
\sum_{\{j,j'\} \in E(H)} (z_{ij} z_{i'j'} + z_{ij'} z_{i'j}) -1,
\]
which is what we sought to derive. The derivation of $(Z^{(2)}_E)^T Z^{(2)}_E = \Id$ is similar.

Finally, we show how to derive the equation $(Z^{(2)}_E)^T D(G) Z = D(H)$ from the equations $ZA(G) = A(H) Z$, where $A(G)$ denotes the adjacency matrix of $G$, with $A(G)_{ii'} = 1$ iff $\{i,i'\} \in E(G)$. Writing out the equations in indices, we need to derive
\[
\sum_{\{i,i'\} \in E(G),k \in V(G)} \left(Z^{(2)}_{E}\right)_{\{i,i'\},\{j,j'\}} D(G)_{\{i,i'\}, k} z_{k \ell} = D(H)_{\{j,j'\},\ell} (\forall \ell \in V(H), \forall \{j,j'\} \in E(H))
\]
Using the fact that $D(G)_{\{i,i'\},k} = \delta_{ik} + \delta_{i'k}$ and the definition of $Z^{(2)}$, this is the same as
\[
\sum_{\{i,i'\} \in E(G),k \in V(G)} \left(z_{ij} z_{i'j'} + z_{ij'} z_{i'j} \right) (\delta_{ik} + \delta_{i'k}) z_{k \ell} = \delta_{j\ell} + \delta_{j'\ell} (\forall \ell \in V(H), \forall \{j,j'\} \in E(H))
\]
Thus we need to derive:
\[
\sum_{\{i,i'\} \in E(G)} \left(z_{ij} z_{i'j'} + z_{ij'} z_{i'j} \right) (z_{i \ell}  + z_{i' \ell})= \begin{cases}
1 & \ell \in \{j,j'\} \\
0 & \text{otherwise}.
\end{cases}
\]
Expanding out the summand, we find the four terms
\[
z_{ij} z_{i'j'} z_{i\ell} + z_{ij} z_{i'j'} z_{i' \ell} + z_{ij'} z_{i'j} z_{i\ell} + z_{ij'} z_{i'j} z_{i'\ell}.
\]
When $\ell \notin \{j,j'\}$, each of these terms is divisible by one of the monomial (local non-isomorphism) axioms, respectively: $z_{ij} z_{i\ell}$, $z_{i'j'} z_{i'\ell}$, $z_{ij'} z_{i\ell}$, and $z_{i'j} z_{i'\ell}$.

Finally, when $\ell \in \{j,j'\}$, without loss of generality suppose that $\ell = j$. Then the only terms that are not divisible by the monomial axioms as above are $z_{ij}^2 z_{i'j'} + z_{ij'} z_{i'j}^2$. Using the Boolean axioms we can easily convert each such summand to $z_{ij} z_{i'j'} + z_{ij'} z_{i'j}$. The derivation of the sum of these over all $\{i,i'\} \in E(G)$ is analogous, \emph{mutatis mutandis}, to the derivation of (\ref{eq:Z2}) above. This completes the proof.
\end{proof}


\begin{proposition} \label{prop:MonTI}
The many-one reduction from \textsc{Monomial Code Equivalence} to \textsc{Tensor Isomorphism} from Grochow \& Qiao \cite{GQ2} is in fact a (2, 4)-many-one reduction.
\end{proposition}

\begin{proof}
We recall the reduction, then prove that it is a low-degree PC reduction. Let $M$ be a $k \times n$ matrix. We build a 3-tensor of size $(k+2n) \times n \times (1 + 2n)$ as follows. The first frontal slice is $\begin{bmatrix} M \\ 0_{2n \times n} \end{bmatrix}$. The remaining $2n$ slices all have just a single nonzero entry, which serve to place a $2 \times 2$ identity matrix ``behind and perpendicular'' to $M$, one $2 \times 2$ matrix in each column. Let us index these slices by $[n] \times 2$. Then the $(i,b)$ slice has a $1$ in entry $(2(i-1) + b, i)$, for all $i \in [n], b \in [2]$. Let us call this tensor $r(M)$. Then the reduction maps $M,M'$ to $r(M), r(M')$.

Let $X,X',Y,Y',Z,Z'$ be the variable matrices for the \textsc{TI} equations for $r(M), r(M')$, and let $A,B,A',B'$ be the variable matrices for \textsc{Monomial Code Equivalence} of $M,M'$ (that is, $AMB^T = M'$, $A$ is invertible, $B$ is monomial and invertible). Consider the substitution:
\[
X \mapsto A \oplus (B \otimes I_2) \qquad Y \mapsto B \qquad Z \mapsto 1 \oplus (B' \circ B') \otimes I_2
\]
\[
X' \mapsto A' \oplus (B' \otimes I_2) \qquad Y' \mapsto B' \qquad Z' \mapsto 1 \oplus (B \circ B) \otimes I_2.
\]
As before, $B \circ B$ denotes the Hadamard or entry-wise product. Let us see what the \textsc{TI} equations become under this substitution. We get
\[
AMB^T = M' \qquad AA' = A'A = \Id \qquad BB' = B'B = \Id \qquad (B' \circ B')(B \circ B) = (B \circ B)(B' \circ B') = \Id
\]
Indeed, notice that the effect of the $B \otimes I_2$ in $X$ and the $B$ in $Y$ is that the row and column locations of the $2 \times 2$ matrix gadgets get permuted in the same way, and the gadget get multiplied by the \emph{square} of the nonzero entries of $B$. These are then multiplied by the $B' \circ B'$ in $Z$.

Now, we derive these equations from the equations for \textsc{Monomial Code Equivalence}. The first three are already present in the equations for \textsc{Monomial Code Equivalence}. The last one we expand out, to see that we need to derive:
\[
\sum_{j} b_{ij}^2 (b'_{jk})^2 = \delta_{ik} (\forall i,k)
\]
Now, for $i \neq k$, we may take the equation $\sum_j b_{ij} b'_{jk}$ and square it, to derive
\[
\sum_{j \neq j'} b_{ij} b'_{jk} + b_{ij'} b'_{j'k} + \sum_j b_{ij}^2 b_{j'k}^2.
\]
Each term in the first sum is divisible by one of the monomial axioms $b_{ij} b_{ij'}$ since $j \neq j'$, and the second sum is what we wanted to derive.

Finally, for $i = k$, we square the equation $\sum_j b_{ij} b'_{ji} - 1$ and add to it $2 \left(\sum_j b_{ij} b'_{ji} -1 \right)$. We then proceed to cancel terms with the monomial axioms as above, and end up with $\sum_j b_{ij}^2 (b'_{ji})^2 - 1$, as desired.
\end{proof}
